\shyph % slovenské delenie slov pre pdfcsplain
\input opmac
\margins/1 a5 (18,18,18,18)mm % nastavenie formátu A5 a všetky okraje na 18mm
\activettchar"   % nastavuje znak " pre zahájenie a ukončenie verbatim textu
\emergencystretch=2em
\hyperlinks{}{}  % nastavenie č/b klikacích odkazov

%% nastaví lámanie riadkov medzi \begtt ... \endtt, viď 
%% https://petr.olsak.net/opmac-tricks.html#verbbreak
\def\tthook{\rightskip=0pt plus1fil
  \def\ {\discretionary{\hbox{$\;\setminus$}}{}{\kern.5em}}%
  \everypar={\countspaces}}
\def\countspaces#1\fi{#1\fi\hangindent=\ttindent \countspacesA}
\def\countspacesA{\futurelet\next\countspacesB}
\def\countspacesB{\ifx\next\spacemacro \expandafter\countspacesC\fi}
\def\countspacesC#1{\advance\hangindent by.5em\ \countspacesA}
\def\spacemacro{\ }

%% Nepoužil sa program vlna pre zdrojový text, preto:
\input  encxvlna



%%% Vlastný text %%%

\nonum\sec Úprava skriptu vimura

Pôvodný skript vimura z \url{https://gist.github.com/vext01/16df5bd\|48019d451e078}, 
používa pdf prehliadač Zathura a textový editor Vim.
Pomocou neho je možné používať obojsmernú synchronizáciu medzi zdrojovým
textom editovaným editorom Vim, ktorý je doplnený o niekoľko makier a
výsledným pdf súborom zobrazeným prehliadačom Zathura, ktorý túto
synchronizáciu podporuje pomocou rozšírenia Synctex. Obmedzením je jeho
použitie len na jednoduché projekty, kedy zdrojový tex súbor a výsledný pdf
súbor musia mať rovnaký názov súboru, rozdielna je len ich prípona. Pre
väčšie projekty zložené z viacerých zdrojových súborov sa preto nehodí.

Pre odstránenie tohto obmedzenia vznikol upravený skript, ktorý používa pre
názov hlavného súboru (bez jeho prílohy je rovnaký ako zobrazovaný pdf
súbor) premennú uloženú v súbore "~/.tlfname" (pre rovnaký účel používa
tento súbor aj skript texloop -- odtiaľ jeho názov) a číslo procesu
príslušného terminálu, kde beží skript vimura v súbore "~/.ttynum" -- tento
slúži na zobrazenie logu z behu \TeX-u pri jeho preklade. Pre ich načítanie
vo Vim-e je potrebné doplniť konfiguračný súbor "~/.vimrc" o tieto riadky:

\ttline=0\begtt
let fname = readfile(glob('~/.tlfname'))[0]
let ttynum = readfile(glob('~/.ttynum'))[0]
\endtt

Tieto súbory vytvára skript vimura, preto je ho vhodné spustiť už pred prvým
spustením Vim-u s takto doplneným súborom ".vimrc". Ten je potrebné doplniť
aj funkciu Synctex:

\begtt
function! Synctex()
   execute "silent !zathura --synctex-forward " . line('.') . ":" . col('.') . ":" . bufname('%') . " " . g:syncpdf 
  redraw! 
endfunction
\endtt

Pre použitie tejto funkcie namapujeme jej spustenie v príkazovom (map) aj vo
vkladacom režime (imap) na kombináciu kláves "Ctrl-Enter" doplnením týchto 
riadkov do ".vimrc":

\begtt
map <C-enter> :call Synctex()<CR>
imap <C-enter> <ESC> :call Synctex()<CR>i
\endtt

Obdobne pre obidva režimy namapujeme pre kombináciu kláves "Ctrl-F5" preklad
zdrojového tex súboru pomocou pdfcsplain do výsledného pdf súboru doplnením
týchto riadkov do ".vimrc":

\begtt
map <C-F5> <ESC>:wa<CR>:exec ":redir! > " . ttynum <CR>:exec "!pdfcsplain -synctex=1 " . fname <CR><CR> :exec ":redir END" <CR>
imap <C-F5> <ESC>:wa<CR>:exec ":redir! > " . ttynum <CR>:exec "!pdfcsplain -synctex=1 " . fname <CR><CR>:exec ":redir END" <CR> i
\endtt
kde sa počas prekladu zobrazuje jeho priebeh v okne samotného vim-u (či
podľa jeho konfigurácie v skripte vimura jeho grafickom variante gvim), kde
sa po preklade znovu objaví editovaný text. Po preklade je možné jeho
priebeh skontrolovať v okne terminálu, odkiaľ bol spustený bash skript vimura.

\bye
